\documentclass[11pt]{article}
\usepackage[T1]{fontenc}
\usepackage{lmodern,microtype}
\usepackage[letterpaper,margin=0.88in]{geometry}
\usepackage{amsmath,amssymb,amsthm,mathtools,booktabs,longtable,array}
\usepackage{xcolor,enumitem,fancyhdr}
\usepackage[breaklinks,colorlinks=true,linkcolor=blue!45!black,urlcolor=blue!45!black,citecolor=blue!45!black]{hyperref}
\definecolor{navy}{RGB}{24,43,68}
\definecolor{teal}{RGB}{18,105,116}
\setlength{\parindent}{0pt}
\setlength{\parskip}{5pt}
\setlength{\headheight}{14pt}
\pagestyle{fancy}
\fancyhf{}
\fancyhead[L]{\small\color{navy}Ten-variable Diophantine consistency}
\fancyhead[R]{\small 23 October 2026}
\fancyfoot[C]{\thepage}
\setlist{nosep,leftmargin=1.5em}
\allowdisplaybreaks[2]
\newtheorem{theorem}{Theorem}[section]
\newtheorem{lemma}[theorem]{Lemma}
\newtheorem{proposition}[theorem]{Proposition}
\newtheorem{corollary}[theorem]{Corollary}
\theoremstyle{remark}
\newtheorem{remark}[theorem]{Remark}
\newcommand{\Z}{\mathbb Z}
\newcommand{\N}{\mathbb N}
\newcommand{\Q}{\mathbb Q}
\newcommand{\Zp}{\mathbb Z_{>0}}
\newcommand{\ZFC}{\mathsf{ZFC}}
\newcommand{\Con}{\operatorname{Con}}
\newcommand{\Pos}{\operatorname{Pos}}
\newcommand{\file}[1]{\path{#1}}
\newcommand{\pin}{\texttt{c1a102936aca437e9999530a0111e6ddab56dc6d}}
\hypersetup{pdftitle={Ten-variable Diophantine consistency: four square roots and small rational residues},pdfauthor={Prepared for Vladimir Reshetnikov; AI-assisted research},pdfsubject={Explicit Diophantine equations equivalent to the consistency of ZFC},pdfkeywords={Diophantine equations, consistency, ZFC, Pell equations, arithmetic circuits}}

\begin{document}
\thispagestyle{empty}
{\sffamily\bfseries\color{teal} RESEARCH REPORT \hfill 23 OCTOBER 2026\par}
\vspace{0.9cm}
{\Huge\bfseries\color{navy} Ten-variable\\[3pt] Diophantine consistency\par}
\vspace{0.5cm}
{\Large Four square roots, small rational residues,\\and a seventeen-integer companion\par}
\vspace{0.6cm}
Prepared for Vladimir Reshetnikov\\
AI-assisted mathematical research with literal arithmetic circuits,\\
fixed integer payloads, ordinary proofs, and independent arithmetic checks.
\vspace{0.7cm}
\begin{abstract}
We construct an explicitly specified polynomial in \textbf{ten strictly
positive integer variables}, of exact degree \textbf{1152}, evaluable in
\textbf{131 binary arithmetic operations}. For the inherited, fixed compact44
coefficient tuple, ordinary mathematics formalizable in ZFC proves that its
positive solvability is equivalent to the existence of a finite ZFC refutation.
Consequently ZFC proves that its unsolvability is equivalent to
$\Con(\ZFC)$.

The main new step repairs a previously unresolved four-root merger.
Instead of assuming that two moduli are coprime, we first recover the Pell
indices and then show that the remaining integral rational residue has
absolute value less than one. It must vanish. The same estimate rejects
the negative main Pell root. Thus one coordinate disappears without imposing
any new divisibility condition on the canonical witnesses.

A separate integer construction merges two roots using adjacent-coefficient
Pell descent. It gives \textbf{seventeen unrestricted integer variables at
degree 112}, together with an eighteen-integer degree-80 companion and
several intermediate tradeoffs. These are improvements relative to the
inspected repository and incoming archives; global minimality and priority
in the wider literature are not asserted.

The ten-variable equation retains the native compiler and K215 endpoint,
but uses neither the Jones nine-variable compression nor Sun's integer
compression. The new reductions are proved here. The inherited native
semantic calibration and the outstanding object-language formalization are
identified explicitly; this report does not claim a new end-to-end Lean or
Rocq certificate.
\end{abstract}
\vfill
\small Reviewed repository snapshot: \pin.\\
The complete equation, fixed coefficient data, and reproducible checks are
included in the accompanying archive.
\newpage
\begingroup
\small
\setlength{\parskip}{0pt}
\tableofcontents
\endgroup
\newpage

\section{Outcome, comparison, and logical scope}

\subsection{The theorem being sought}
Fix the recursive syntactic ZFC refutation predicate calibrated by the
repository's native source. Its consistency statement is the arithmetical
sentence asserting that no finite refutation exists. The target is a fixed
integer polynomial $P$, with a stated witness domain, such that
\begin{equation}\label{eq:logical-target}
 \ZFC\vdash
 \bigl[\neg\exists\mathbf x\in\Zp^{10}\ P(\mathbf x)=0\bigr]
 \longleftrightarrow \Con(\ZFC).
\end{equation}
We prove a solvability/refutation equivalence and negate it. Neither
consistency nor inconsistency is assumed. A proof of the displayed
equivalence does not separately prove either side.

The new argument is an ordinary mathematical proof, using integer
arithmetic, finite sequences, Pell recurrences, and finite algebraic
extensions. These arguments can be carried out in ZFC. The repository's
separate effort to produce an explicit derivation in its encoded ZFC
calculus remains a different verification task. The snapshot's
\file{STATUS.md} and \file{OBJECT_ZFC.md}, retained in the archive, describe
that boundary \cite{active}.

\subsection{Measured results}
The following table compares specific constructions, not a hypothetical
single ordering of all equations by simplicity. An operation is one live
binary addition, subtraction, or multiplication; squaring and multiplication
by a fixed scalar are charged. Fixed numeral leaves and reuse are free.

\begin{center}
\small
\begin{tabular}{@{}lrrrr@{}}
\toprule
Construction & Domain & Variables & Degree & Operations\\
\midrule
Active shared-product baseline & $\Zp$ & 9 & 1,845,940 & 440\\
Package 71 lower-degree baseline & $\Zp$ & 9 & 1,379,730 & 478\\
Incoming radical-guard reduction & $\Zp$ & 9 & 1,158,228 & 1753\\
Incoming centered three-root form & $\Zp$ & 11 & 257 & 114\\
\textbf{Four-root residue equation} & $\mathbf{\Zp}$ & \textbf{10} & \textbf{1152} & \textbf{131}\\
Translated companion & $\N$ & 10 & 1152 & 141\\
\midrule
Incoming joint-guard endpoint & $\Z$ & 19 & 88 & 109\\
\textbf{Centered pivot companion} & $\mathbf\Z$ & \textbf{18} & \textbf{80} & 115\\
\textbf{Two-root integer merger} & $\mathbf\Z$ & \textbf{17} & \textbf{112} & 124\\
Sheared integer companion & $\Z$ & 17 & 110 & 140\\
\bottomrule
\end{tabular}
\end{center}
The ten-positive-variable result occupies a useful gap: one more variable
than the nine-variable constructions, but much smaller degree; one fewer
variable than the direct eleven-variable construction, at a moderate
increase in degree. Its graph contains 77 multiplications, 37 additions,
and 17 subtractions. The natural companion replaces every positive
coordinate $x$ by $u+1$, adding ten gates and preserving degree.

The integer companion uses a different fixed tuple and is compared only
with integer-domain constructions. Ten unrestricted integers are already
available through Sun's compression \cite{sun,structural}; we do not claim
seventeen is the fewest possible integer variables. The point is the
degree/variable tradeoff, without that compression. The supplied integer graphs have measured operation counts; no global circuit minimum is claimed.

\subsection{What was actually reviewed}
We inspected the placed packages 31--75 and all five incoming archives at
the pinned snapshot. Two of the incoming packages already combine signed
power recovery, determinant guards, and coordinate projections. Another
already proves a four-radical separation lemma but leaves carrier
compatibility open. The current report builds on those facts rather than
claiming the separation lemma or a previously delivered pivot as new.

The small incoming \file{report.zip} contains only a main TeX file with
missing inputs. Its announced numerical results are distinguished from
complete, independently replayable packets in our audit. The four larger
archives provide the actual comparison data \cite{joint,structural,radical,budget}.
Our archive records their Git blob identities and SHA-256 hashes.

\section{The complete ten-variable polynomial}\label{sec:equation}

\subsection{Domain and fixed coefficients}
The ten independent variables are exactly
\begin{equation}\label{eq:variables}
 j,F,\alpha,i,T,X,s,n,\eta,\zeta\in\Zp.
\end{equation}
Every other symbol in this section abbreviates an integer polynomial in
these variables. No abbreviation introduces another quantified variable.

The four fixed coefficients are the coherent compact44 tuple
\[
 (\beta,K_0,M_C,M_F),\qquad \beta=2^{590069}-1.
\]
The high mask $M_F$ is already shifted by $\beta$; it must not be shifted
again. The archive contains all four actual integers as compressed,
unsigned little-endian byte strings, and reconstructs them from the
literal 430-action K215 table using the pinned compact44 recipe. Thus the
ports are fixed integer coefficients, not free parameters.

The coefficient bounds used below are
\begin{equation}\label{eq:maskbounds}
 K_0\ge1,\quad 0<M_C<\beta,\quad 0<M_F<2\beta,
 \quad \beta\equiv1\pmod3.
\end{equation}
Their full source meaning is supplied by the compact44 compiler adapter
\cite{compact44}. The arithmetic size of the four payloads is given in
Appendix~\ref{app:coefficients}.

\subsection{Outer and Pell registers}
Define
\begin{align*}
 J&=3j+1,&q&=\beta J+1,&C&=q-F-\alpha,\\
 Y&=q^3s,&E&=XY,&k&=\eta+\zeta,\\
 c&=kY+\eta,&a&=E+Y,&H&=4a+3,\\
 \Delta&=a^2+H=(a+2)^2-1,&\ell&=EkY.
\end{align*}
Next set
\begin{align*}
 R&=\bigl((q-1)(q-F)+\alpha\bigr)(q^2-1)+(M_C+qM_F)J,\\
 \widehat R&=q^3(q+F+\alpha),&
 B_t&=(qK_0+X)C-qF,\\
 I&=ic^2,&D_f&=1+\Delta I^2,\\
 Q&=\Delta^2I^2,&\delta&=Q-1,\\
 V&=D_f(cT-R)-c,&D_t&=\ell(\ell+k)+1,\\
 D_d&=1+\Delta c^2,&\overline V&=D_f(cT+\widehat R),\\
 \overline\Gamma&=2c+2X+k+\widehat R+1,&
 L&=3\overline V+\overline\Gamma,\\
 W&=1+(q-1)L,&v&=k-R-1,\\
 U&=q^5,&b_0&=X+ac,\qquad M=(q-1)U.
\end{align*}
The combined carrier is the polynomial
\begin{equation}\label{eq:Z}
 Z=M\bigl[H(\ell+k+n+L+s)+b_0\bigr]
       -HUWB_t+H(q-1)(q^4X-v).
\end{equation}
Its four scaled radicands are represented without division by
\begin{equation}\label{eq:radicands}
 A=M^2H^2D_t,\quad B=M^2H^2D_f,\quad
 D=M^2D_d,\quad C_*=M^2H^2(QV^2-1).
\end{equation}
Here $C_*/\delta$ will occur only in the proof, not in the equation.

\subsection{A short division-free norm}
Define the final registers
\begin{align*}
 \sigma&=A+B+D,&\pi&=AB+(A+B)D,&\omega&=ABD,\\
 a_*&=\delta(Z^2-\sigma)+C_*,&t_*&=\delta C_*Z^2,\\
 k_*&=a_*^2+4(t_*-\delta^2\pi),&r_*&=16\delta^3\omega,\\
 b_*&=a_*k_*-r_*,&&\\[-5pt]
 e_*&=k_*^2+16t_*a_*^2-4r_*(\delta Z^2+C_*).&&
\end{align*}
The equation is
\begin{equation}\label{eq:P10}
 \boxed{P_{10}=e_*^2-64t_*b_*^2=0.}
\end{equation}
All powers are fixed nonnegative integer powers. This is an ordinary
polynomial equation with integer coefficients, not an exponential
Diophantine equation, a bounded-quantifier abbreviation, or an oracle call.
Its literal graph is \file{ten_positive/generated/p10_positive.json}.

\begin{theorem}\label{thm:main}
For the four fixed coefficients specified above,
\[
 \exists\mathbf x\in\Zp^{10}\ P_{10}(\mathbf x)=0
 \quad\Longleftrightarrow\quad
 \text{there is a finite syntactic ZFC refutation}.
\]
This equivalence has an ordinary proof formalizable in ZFC. In particular,
\eqref{eq:logical-target} holds. The polynomial has exact total degree
1152, and the delivered arithmetic graph has 131 live binary operations.
\end{theorem}

\section{What the four-root norm means}

For rational arguments define the ordinary algebraic norm polynomial
\[
 F_{16}(z;a,b,d,c)=
 \prod_{\varepsilon\in\{-1,1\}^4}
 (z-\varepsilon_1\sqrt a-\varepsilon_2\sqrt b
     -\varepsilon_3\sqrt d-\varepsilon_4\sqrt c).
\]
This symmetric product is a polynomial in $z,a,b,d,c$ with integer
coefficients. A useful three-root identity is
\begin{equation}\label{eq:F8}
 F_8(z;a,b,d)=
 \bigl[(z^2-a-b-d)^2-4(ab+ad+bd)\bigr]^2-64abdz^2.
\end{equation}
Expanding $F_8(z+\sqrt c)F_8(z-\sqrt c)$ into its even and odd parts
gives the final schedule in Section~\ref{sec:equation}. More precisely,
\begin{equation}\label{eq:normidentity}
 P_{10}=\delta^8F_{16}(Z;A,B,D,C_*/\delta).
\end{equation}
The identity is polynomial after clearing denominators, and is checked
by exact arithmetic independently of the graph builder.

At every positive input, $\delta,H,M>0$, and all four radicands are
positive. Hence a zero of $P_{10}$ supplies signs such that
\begin{equation}\label{eq:root-equality}
 \frac{Z}{MH}=\tau_*+f_*+y_*+\frac{d_*}{H},
 \qquad
 \begin{cases}
 \tau_*^2=D_t,\\
 f_*^2=D_f,\\
 y_*^2=(QV^2-1)/(Q-1),\\
 d_*^2=D_d.
 \end{cases}
\end{equation}
We must prove these roots are integral. Their integrality does not follow
merely from taking a norm.

\section{Two lemmas for recovering the roots}

\subsection{Every nonempty subset must be separated}
\begin{lemma}\label{lem:rational}
Let $r_1,\ldots,r_m$ be positive real square roots of positive rational
numbers. Suppose every nonempty signed subset sum of the $r_i$ is
nonzero. If a signed sum of all $r_i$ is rational, every $r_i$ is rational.
\end{lemma}
\begin{proof}
In the finite multiquadratic field generated by these roots, group the
terms according to their square classes in $\Q^*/\Q^{*2}$. Distinct
square classes are linearly independent over $\Q$: averaging a proposed
relation against the corresponding quadratic characters of the finite
Galois group isolates each class. In a rational signed sum, every
nontrivial square-class contribution must therefore vanish. A nonempty
such class would be a forbidden signed subset sum. All classes are trivial.
\end{proof}

For four roots, distinct radicands and a nonzero total signed sum are
insufficient. The identity
\[
 \sqrt2+\sqrt8-\sqrt{18}+100=100
\]
contains a cancelling triple. One globally dominant root does not repair
this defect. The stronger hypothesis in Lemma~\ref{lem:rational} is the
one we actually establish.

Let $t,d,f,y$ denote the positive roots of $D_t,D_d,D_f,(QV^2-1)/(Q-1)$.
Positivity of the ten inputs alone gives
\begin{align*}
 q&\ge16,&Y&\ge q^3\ge4096,&k&\ge2,&c&\ge2Y+1,\\
 \ell&<t<\ell+k<(a+1)c,&f&>(a+1)c^2,&0&<d/H<c.
\end{align*}
Since $c<k(Y+1)$ and $\ell=XkY^2$, we have $t>2c+1$.
Write $b=cT-R\in\Z$. If $b=0$, then $V=-c$, so
$c<y<c+1$; therefore the ordering $d/H,y,t,f$ has each term after the
first strictly greater than the sum of its predecessors. If $b\ne0$,
then $|V|\ge D_f-c=f^2-c>2f$, and $y>|V|$. The ordering
$d/H,t,f,y$ has the same property. Here $f-t>(a+1)c(c-1)$ and
$f>3c^2$ supply the stated strict gaps.

Thus every nonempty signed subset is separated before any source
equation, sign of $C$, or divisibility conclusion is used.
Lemma~\ref{lem:rational} applied to \eqref{eq:root-equality} makes
$\tau_*,f_*,d_*/H,y_*$ rational. The first three unscaled roots are
integers because their squares are integers.

\subsection{Adjacent-coefficient Pell descent}
\begin{lemma}\label{lem:descent}
If $Q>1$ and $V$ are integers, $y\in\Q$, and
\[
 QV^2-(Q-1)y^2=1,
\]
then $y\in\Z$.
\end{lemma}
\begin{proof}
Use absolute values and write $y=a/m>0$ in lowest terms, with $m>0$.
Then $m^2\mid Q-1$. Neither coordinate is zero. If $V=1$, then $y=1$.
For $V>1$ define
\[
 V'=(2Q-1)V-2(Q-1)y,\qquad
 y'=(2Q-1)y-2QV.
\]
The first coordinate is integral because $m\mid Q-1$. The norm is
unchanged. The inequalities $y>V$ and $QV>(Q-1)y$ imply
$-V<V'<V$, while the norm excludes $V'=0$. The reduced denominator
of $y'$ is still $m$, since $2Q-1\equiv1\pmod m$.
Strong induction on $|V|$ reaches $V=1$ without changing that denominator;
hence $m=1$.
\end{proof}
This lemma, used in the incoming Pell constructions \cite{cones,joint},
now gives $y_*\in\Z$. No coprimality of another modulus with $Q-1$ is
needed. Adjacent coefficients matter: $2\cdot1^2-4(1/2)^2=1$ is a
counterexample for a general pair of coefficients.

\section{Positivity before exact divisibility}\label{sec:bootstrap}

The fixed mask bounds imply, for every positive input,
\begin{equation}\label{eq:majorants}
 |R|<\widehat R,\qquad |V|<\overline V,\qquad f+y<3\overline V.
\end{equation}
For the first bound, the mask contribution is less than
$(q-1)(1+2q)<2q^2$, and comparison with $q^3(q+F+\alpha)$ is strict
for $q\ge16,\alpha\ge1$. The gap $\widehat R-|R|\ge1$, together
with $D_f>c$, gives $|V|<D_f(cT+\widehat R)$. Finally
$f<D_f<\overline V$ and $y<2|V|<2\overline V$.

Define two proof-only rational coordinates
\begin{align}
 r&=\tau_*-\ell-k-s+\frac{WB_t}{q-1},\label{eq:rproof}\\
 \gamma&=\frac{d_*-b_0}{H}-\frac{q^4X-v}{q^5}.\label{eq:gproof}
\end{align}
For either sign of $d_*$,
\[
 |\gamma|
 \le\frac dH+\frac{b_0}H+\frac Xq+\frac{|v|}{q^5}
 <2c+2X+k+\widehat R+1=\overline\Gamma.
\]
Equation~\eqref{eq:root-equality} is exactly
\begin{equation}\label{eq:merge}
 r+\gamma+f_*+y_*=n+L.
\end{equation}
Thus
\[
 r\ge n+L-f-y-|\gamma|>n>0.
\]
At this stage $r$ and $\gamma$ may still be rational. The strict
inequality is enough. Since $|\tau_*|<\ell+k$,
\[
 \frac{WB_t}{q-1}=\ell+k+s+r-\tau_*>s+r>0.
\]
Therefore $B_t>0$, which forces $C>0$. The authenticated outer packing
bounds now give
\begin{equation}\label{eq:Rbounds}
 3q<R<q^4-2.
\end{equation}
This order is central: positivity is recovered over the rationals before
the two congruence channels are separated over the integers.

\section{Recover the Pell indices without choosing the main sign}

For $b\ge2$ define the Pell polynomials by
\[
 \chi_b(m)+\psi_b(m)\sqrt{b^2-1}
 =(b+\sqrt{b^2-1})^m.
\]
Put $A_0=a+2$ and $P_0=2XY^2+1$. Positive Pell classification gives
\begin{equation}\label{eq:classify}
 |d_*|=\chi_{A_0}(p),\quad c=\psi_{A_0}(p),\qquad
 |\tau_*|=\chi_{P_0}(m),\quad k=2\psi_{P_0}(m).
\end{equation}
Only the positive ordinate $c$ and absolute roots occur here. No main-root
orientation or positive $\gamma$ is required.

Since $c>1$, $p\ge2$. If $m=1$, then $k=2$ and
$\eta=\zeta=1$, giving $c=2Y+1<2A_0=\psi_{A_0}(2)$, a contradiction.
Hence $m\ge2$. Since $P_0>A_0$ and $c>k$, monotonicity gives $m<p$.
Moreover
\begin{equation}\label{eq:cgrowth}
 c>kY\ge4P_0Y>8XY^3\ge8q^9>2R,
 \qquad c>2p.
\end{equation}

The normalized strong equation is already encoded by $D_f$:
\[
 |f_*|^2-\Delta(ic^2)^2=1.
\]
Its Pell index $m_s$ satisfies $pc\mid m_s$. Indeed, strong divisibility
of the $\psi_{A_0}$ sequence gives $p\mid m_s$. Writing $m_s=pv_s$
and reducing the quotient identity
\[
 \psi_{A_0}(pv_s)/\psi_{A_0}(p)=\psi_{\chi_{A_0}(p)}(v_s)
\]
modulo $c$ shows $c\mid v_s$, since
$\chi_{A_0}(p)^2\equiv1\pmod c$. In particular $c\mid m_s$ and
$m_s>2p$.

For clarity, the inherited signed auxiliary-index lemma used next has
the following hypotheses: $c=\psi_{A_0}(p)$,
$f=\chi_{A_0}(m_s)$, $c\mid m_s$, $m_s>2p$;
$S\ge2$, $S^2=\Delta(f^2-1)$, and $c\mid S$;
$S^2V^2-(S^2-1)y^2=1$;
$V\equiv-c\pmod f$, $V\equiv-R\pmod c$; and
$0<p,R<c/2$. It permits signed $V,y$ and does not assume any main-root
congruence \cite[lemmas on odd subsequences, signed index recovery, and
normalized rank]{carriers}.

Here take $S=\Delta I$. Then $S^2=Q$, while
$V=D_f(cT-R)-c$ has both required congruences because $f^2\equiv1\pmod c$.
All interval hypotheses follow from \eqref{eq:cgrowth}. The lemma gives
\begin{equation}\label{eq:pR}
 p=R.
\end{equation}
Its proof uses the odd Pell identities to obtain
$R\equiv\pm p\pmod c$, followed by the strict half-interval bounds;
the full inherited argument is included in the archive.

\section{Separate the two residues by size}\label{sec:residue}

\subsection{The q-power channel}
All starred roots are now integral, so \eqref{eq:merge} implies
$r+\gamma\in\Z$. Clearing the denominators in
\eqref{eq:rproof}--\eqref{eq:gproof} gives
\[
 q^5\mid H(q-1)(q^4X-v).
\]
Since $q\equiv2\pmod3$, $q^3\mid Y$, and $H=4Y(X+1)+3$,
we have $\gcd(H,q)=1$. Also $\gcd(q,q-1)=1$. Consequently
\begin{equation}\label{eq:q5}
 q^5\mid q^4X-v,
 \qquad q^4\mid v.
\end{equation}
The first Pell recurrence gives $k\equiv2m\pmod{4XY^2}$, whose modulus
is divisible by $q^4$. Thus $2m-R-1\equiv0\pmod{q^4}$.
Since $1\le m<p=R<q^4$, that difference has absolute value less than
$q^4$. It is zero:
\begin{equation}\label{eq:firstindex}
 2m=R+1.
\end{equation}
Now $v=k-2m$ is a positive multiple of $4XY^2$, hence of $q^5$.
Equation~\eqref{eq:q5} therefore forces $q\mid X$.
Returning to the merged equality, whose $q^5$ fraction is now integral,
gives
\begin{equation}\label{eq:weakH}
 \frac{(q-1)(d_*-X-ac)}H\in\Z.
\end{equation}
We have not cancelled $q-1$ modulo $H$.

\subsection{The bound which makes cancellation unnecessary}
Put $r_0=(R-1)/2=m-1\ge24$. Recurrence induction gives
\[
 \psi_b(t)>(2b-1)^{t-1}\quad(t\ge2),\qquad
 \psi_b(t)\le(2b)^{t-1}\quad(t\ge1).
\]
The exact identity
\[
 (2A_0-1)^2-2P_0X
 =4Y^2(2X+1)+12Y(X+1)+9-2X>0
\]
yields
\[
 \frac ck=\frac{\psi_{A_0}(2r_0+1)}{2\psi_{P_0}(r_0+1)}
 >\frac{X^{r_0}}2.
\]
But $c/k=Y+\eta/k<Y+1$, so, using $X\ge q\ge16$,
\begin{equation}\label{eq:Yhuge}
 Y>\frac{X^{r_0}}2-1\ge\frac{q^{r_0}}2-1
 >q\,2^{2r_0+1}=q\,2^R.
\end{equation}
The last elementary inequality already holds for $r_0\ge4$.

\subsection{Reject the negative main root and recover the power}
The Pell polynomial identities modulo $H$ are
\begin{align}
 \chi_{A_0}(R)-a\psi_{A_0}(R)&\equiv2^R\pmod H,\label{eq:pellmod1}\\
 2^R\bigl(\chi_{A_0}(R)+a\psi_{A_0}(R)\bigr)&\equiv1\pmod H.
 \label{eq:pellmod2}
\end{align}
They follow by substituting $\Delta\equiv a^2\pmod H$ into the Pell
binomial expressions: $A_0-a=2$, $A_0+a=2a+2$, and
$2(2a+2)=H+1$.

If $d_*=-\chi_{A_0}(R)$, then \eqref{eq:weakH} and
\eqref{eq:pellmod2} imply
\[
 \frac{(q-1)(1+2^RX)}H\in\Z.
\]
This number is positive but, by \eqref{eq:Yhuge}, is strictly less than
one. That is impossible. Hence $d_*>0$.

For the positive root, \eqref{eq:weakH} and \eqref{eq:pellmod1} give
\[
 \frac{(q-1)(2^R-X)}H\in\Z.
\]
Its absolute value is less than one because
\[
 H>4q\,2^R(X+1)>q(2^R+X)>(q-1)|2^R-X|.
\]
It must be zero. We have proved
\begin{equation}\label{eq:exactpower}
 X=2^R,\qquad H\mid d_*-X-ac.
\end{equation}
The proof-only $\gamma$ and then $r$ are integral. Since
$W\equiv1\pmod{q-1}$, their defining identities also give
\begin{equation}\label{eq:transportdiv}
 q-1\mid B_t.
\end{equation}
This is the new residue-separation argument. It uses an exact integer
interval, with no new condition on $\gcd(H,q-1)$ or on canonical $Y$.

\section{Soundness and completeness of the new equation}

\subsection{Reconstruct a positive normalized source solution}
Define
\[
 w=X/q,\quad h=v/(XY),\quad
 \gamma_{\rm old}=(d_*-X-ac)/H,
 \quad f_{\rm old}=|f_*|,\quad y_{\rm old}=|y_*|.
\]
These are integers. They are positive: $w>0$; $h>0$ follows from
$v=k-2m>0$; and $d_*>(a+1)c$, $c>X$ give
$H\gamma_{\rm old}>c-X>0$.
The first, main, strong, and auxiliary Pell equations now hold, as does
$k-hXY-R=1$.

Because $B_t=q[(K_0+w)C-F]>0$, \eqref{eq:transportdiv} gives the
positive integer
\[
 z_{\rm new}=\frac{(K_0+w)C-F}{q-1}.
\]
Set $z=z_{\rm new}+1$. Then the source transport identity is exactly
\[
 (K_0+w)C+(q-F)-z(q-1)=1.
\]
Finally set $T_{\rm old}=f_{\rm old}T$. The source auxiliary coordinate
agrees because
\[
 c(T_{\rm old}f_{\rm old}-1)-Rf_{\rm old}^2
 =D_f(cT-R)-c=V.
\]
Thus we have a positive solution of the normalized compact44 source from
package 61. Its compiler/calibration theorem gives a finite ZFC refutation.
No search for a refutation is performed by the verification programs.

\subsection{Construct a new zero from a canonical source solution}
Conversely a finite refutation gives canonical source witnesses. The
source already has $Y=q^3s$ with $s>0$; no additional congruence on $s$
is imposed. Its padded repunit $J=1+B+\cdots+B^{N-1}$ has odd $N>1$
and $B=2^{590069}\equiv2\pmod3$. Hence $J\equiv1\pmod3$ and
$j=(J-1)/3>0$.

Choose the normalized strong pair $f,I=ic^2$ and the positive auxiliary
pair as in the source, using its canonical $R\equiv3\pmod4$ condition.
Compute $L,W$ after these choices. The source
gives $t_0=B_t/(q-1)\ge q$, $\ell<\tau<\ell+k$, and $s<k$.
Set
\[
 r=Wt_0+\tau-\ell-k-s,
 \qquad
 \gamma=\gamma_{\rm old}-X/q+v/q^5.
\]
Both are integers; $v$ is a positive multiple of $4XY^2$ and $q\mid X$.
The unconditional bound $|\gamma|<\overline\Gamma<L$ applies.
Also $L>k$, so
\[
 r>qW-2k>2L.
\]
Now take
\[
 n=r+\gamma+f+y-L>r-|\gamma|-L>0.
\]
The four-root equality holds with the positive roots, and therefore
$P_{10}=0$. All ten coordinates are strictly positive. This completes
both directions of Theorem~\ref{thm:main}, subject exactly to the stated
inherited source theorem.

\section{Exact degree and arithmetic cost}

Every input in \eqref{eq:variables} has degree one; the four fixed ports
have degree zero. The relevant register degrees are
\[
\begin{array}{c|rrrrrrrrrrrr}
 &Y&c&\Delta&I&D_f&Q&V&L&Z&A&B&C_*\\\hline
\deg&4&5&10&11&32&42&38&38&51&42&54&140
\end{array}
\]
and $\deg D=32$. In the last norm schedule,
\[
 \deg a_*=144,\quad \deg k_*=288,\quad
 \deg b_*=432,\quad \deg e_*=576.
\]
The subtracted term $64t_*b_*^2$ has degree 1148, below the degree 1152
of $e_*^2$. Thus the leading form is $Q_0^8Z_0^{16}$.

To prove that the degree remains exact on the actual coefficient slice,
put
\[
 q_0=3\beta j,\quad k_0=\eta+\zeta,\quad Y_0=q_0^3s,
 \quad B_{t,0}=(K_0q_0+X)(q_0-F-\alpha)-q_0F.
\]
Then
\begin{align*}
 Q_0&=i^2k_0^4X^4Y_0^8,\\
 Z_0&=-12q_0^6i^2k_0^5X^3Y_0^8T B_{t,0},\\
 [P_{10}]_{1152}
 &=12^{16}q_0^{672}i^{48}k_0^{112}X^{80}s^{192}T^{16}B_{t,0}^{16}.
\end{align*}
In particular the coefficient of
\begin{equation}\label{eq:monomial}
 j^{672}i^{48}\eta^{112}X^{96}s^{192}T^{16}\alpha^{16}
\end{equation}
is exactly
\[
 12^{16}(3\beta)^{672}>0.
\]
The minimum $j$-degree and the absence of $F$ in the target monomial
force $-X\alpha$ from each factor of $B_{t,0}$;
no specialization of $K_0,M_C,M_F$ can cancel this coefficient. The
monomial has degree 1152, proving exactness rather than merely a syntactic
upper bound.

The graph builder shares $q^2,q^3,q^4$, $D_f-1$, $kY$, and the repeated
norm products. The independent checker traverses the live graph, checks
the complete declared input set, and recounts every operation. The result
is $77+37+17=131$. These figures do not measure bit complexity: the
fixed coefficients and potential witnesses are large.

\section{The unrestricted-integer companions}\label{sec:integer}

The following separate reduction starts from the incoming joint-guard
integer family \cite{joint,structural}. It uses the source34 coefficient
tuple, together with two fixed cone scales. It does not silently substitute
the compact44 tuple used in $P_{10}$.

\subsection{A transport pivot beyond the earlier projections}
The retained transport equation becomes monic in $F$ after using the
complement coordinate $C$:
\[
 \alpha=q-F-C,\qquad
 F=(K_0+W)C-z(q-1).
\]
These are polynomial substitutions over $\Z$. Forgetting the eliminated
coordinate and restoring it by this formula are inverse on the integer
zero sets. Applying this previously known pivot to the more deeply
projected joint-guard packets gives the following new combinations:
\[
\begin{array}{c|rrrr}
\text{integer variables}&21&20&19&18\\\hline
\text{exact degree}&16&20&32&88
\end{array}
\]
The pivot raises the packing expression's degree only to five and the
determinant guard's degree only to four, below the existing dominant rows.
The delivered packets include their complete restoration maps and exact
degree certificates.

\subsection{Centering one Pell root without a parity witness}
For the first Pell residual write $u=E(c-\eta)$. Instead of the old root
$\tau_{\rm old}$, set $t=2\tau_{\rm old}-2u-k$. Its residual is
\begin{equation}\label{eq:center}
 t^2+2(2u+k)t+k^2-4=0.
\end{equation}
Modulo four this forces $t+k$ even, so
$\tau_{\rm old}=(t+2u+k)/2$ is an integer automatically. Equation
\eqref{eq:center} is exactly four times the old residual under this
substitution. The high-degree $u^2$ term disappears.

In the eighteen-variable packet this lowers the first residual's degree
from 34 to 18. Let $G$ be the remaining joint positivity residual of
degree 44, and let $S$ be the sum of squares of the other residuals, of
degree 36. Since $S\ge0$ over $\Z$,
\[
 (1+G)(1+S)-1=0\quad\Longleftrightarrow\quad G=0\text{ and }S=0.
\]
Indeed $1+S$ is a positive integer divisor of one. The resulting single
polynomial has exact degree $44+36=80$ in eighteen integers.

\subsection{A second root merger gives seventeen integers}
Retain the strong equation and put
\[
 I=(2r+1)c^2,\quad Q=I^2,\quad
 A=u(u+k)+1,\quad B=\frac{QV^2-1}{Q-1}.
\]
The old first and auxiliary roots satisfy $\tau^2=A$ and $y^2=B$.
Replace their two coordinates by the single integer
\[
 \xi=2(I\tau+y)-I(2u+k).
\]
Write $s_I=(\xi+I(2u+k))/2$ in the proof only. The denominator-cleared
two-root norm is
\begin{equation}\label{eq:integernorm}
 \bigl[(Q-1)(4s_I^2-4QA)-4(QV^2-1)\bigr]^2
 -64(Q-1)QA(QV^2-1)=0.
\end{equation}
Here $4s_I^2=(\xi+I(2u+k))^2$, so the displayed equation is polynomial
in the retained integer variables. The literal packet uses the equivalent
centered expression to avoid unnecessary high-degree cancellation.

The retained joint guard proves the nondegeneracy needed for this merger:
$|c|>1$, $Q>1$, and $A>0$, before the deleted root conditions are used.
The two rational radicands $QA$ and $B$ cannot coincide. Otherwise
\[
 (Q-1)QA=QV^2-1
\]
would contradict reduction modulo $Q$. If $B<0$, the norm cannot have
a real zero; the case $B=0$ is also impossible modulo $Q$. For positive,
distinct radicands, a rational signed sum of the two roots makes both
roots rational. Therefore $\tau$ is integral by rational-square
integrality, and $y$ is integral by Lemma~\ref{lem:descent}. This recovers
the two old rows without an extra parity or divisibility witness.

Conversely, old integer roots give the displayed integer $\xi$. Thus the
root replacement preserves existential solvability in both directions.
Together with the earlier transport pivot this gives a single sum of
squares in \textbf{seventeen unrestricted integers}, of exact degree
\textbf{112}. The archive includes all residuals, their variable list,
coefficient binding, restoration maps, and independent checks. Its supplied graph uses 124 operations, without a claim of operation minimality.

\subsection{Complete display of the seventeen-variable equation}

For an explicit standalone specification, the seventeen unrestricted
integer coordinates are
\[
 J_3,C,z,f,h,r,v,w_3,\gamma,o,j,c,t,p_1,p_2,p_3,\xi.
\]
Here $z$ remains the transport coordinate and $t$ is the joint-guard
root coordinate. Set $\lambda=2^{1113}$ and $\mu=2^{666104}$; the other
four fixed coefficients are the complete source34 tuple
$\beta,K_0,M_C,M_F$, with $\beta=2^{1064921}-1$ and $M_F$ already shifted.
Use only the following polynomial abbreviations:
\begin{align*}
 J&=3J_3+1,&q&=\beta J+1,&W&=4(3w_3+1),\\
 X&=qW,&Y&=(2v+1)q^3,&E&=XY,\\
 a&=(X+1)Y,&H&=4a+3,&\Delta&=a^2+H,\\
 V&=fo-c,&k&=1-V+jc+hE,&\eta&=c-Yk,\\
 u&=EYk,&A&=u(u+k)+1,&I&=(2r+1)c^2,\\
 Q&=I^2,&\delta&=Q-1,&F&=(K_0+W)C-z(q-1).
\end{align*}
Further put
\begin{align*}
 R&=\bigl(q(q-F)-C\bigr)(q^2-1)+(M_C+qM_F)J,\\
 d&=(\lambda F)(\mu C)-(q-\lambda F-\mu C)^2,\\
 T&=\eta(k-\eta),\\
 P_{\rm pos}&=1+p_1^2+p_2^2+p_3(p_3+1),\\
 G&=dT-(t^2-d-T+1)^2-P_{\rm pos},\\
 B_0&=QV^2-1,\\
 L&=\xi^2+2I(2u+k)\xi+Q(k^2-4),\\
 D_n&=\delta L-4B_0,\\
 N&=D_n^2-64\delta QAB_0.
\end{align*}
The promised single equation is
\begin{equation}\label{eq:full17}
 \boxed{P_{17}=M^2+B_p^2+S_p^2+G^2+N^2=0,}
\end{equation}
where
\[
 M=(X+ac+\gamma H)^2-\Delta c^2-1,\quad
 B_p=cj-V-R,\quad S_p=Q-\Delta(f^2-1).
\]
There are no hidden quantified coordinates or divisions in this display.
Its five residual degrees are $14,5,14,44,56$. The coefficient of
\[
 r^{12}c^{24}h^4f^4o^4J_3^{44}w_3^8v^{12}
\]
in $P_{17}$ is exactly $2^{52}3^{52}\beta^{44}>0$, certifying degree
$112$ after specialization to the actual coefficient tuple. The supplied
literal arithmetic DAG uses $124$ binary operations under the stated
counting convention; this is a verified construction, with no claim of
operation minimality.

The scale $I$ in the merger is essential. Taking $I=2$, $Q=4$, $V=2$,
and $A=5$ makes the unscaled auxiliary radicand
$(QV^2-1)/(Q-1)=5$. The unscaled norm then accepts the spurious rational
sum $\sqrt5-\sqrt5=0$, although neither square root is integral. The
scaled radicands are $20$ and $5$, and the same sum-zero input is rejected.

\subsection{A degree-110 integer companion}
A polynomial change of coordinate $\xi_{\rm old}=\xi+2V$ gives a separate
seventeen-integer equation of degree 110. Keep the definition of $L$ above,
now using the new coordinate $\xi$, and set
\begin{align*}
 J_*&=\delta(L+4V\xi)-4(V^2-1),\\
 H_*&=4\delta I(2u+k)V,\\
 N_{110}&=J_*^2+2J_*H_*+64\delta QA(1-V^2)
                  +16\delta^2Q(k^2-4)V^2.
\end{align*}
The exact identity $N_{110}(\xi)=N(\xi+2V)$ follows from
$(2u+k)^2=4A+k^2-4$ and $\delta=Q-1$. Replace $N^2$ by $N_{110}^2$
in \eqref{eq:full17}; the other four rows are unchanged. The coordinate
change is invertible over the integers, so solvability is preserved.
The coefficient of
\[
 r^{12}c^{24}h^4 f^2o^2\xi^2J_3^{44}w_3^8v^{12}
\]
is $2^{52}3^{52}\beta^{44}>0$, proving exact degree 110.
The supplied circuit uses 140 operations. The degree-112 version remains
the principal integer display because its circuit is shorter, at 124
operations; neither choice dominates the other on both measures.

\section{Why the new argument is different}

The new ten-variable merger has three logically separate stages.
First, a strict hierarchy of all four root magnitudes makes a rational
signed sum force rationality of every root. Second, a sufficiently large
positive transport margin establishes $C>0$ while its recovered quotient
is still only rational. Third, Pell growth makes the remaining integral
rational residue too small to be nonzero.

This last stage removes the earlier obstruction. From
$H\mid(q-1)(d-X-ac)$ one generally cannot cancel $q-1$. Nor can
one freely impose $(q-1)\mid Y$: $Y$ is tied to a half-binomial value,
not an arbitrary large slack. The new proof requires neither move.

There is a concrete control even within the exact half-binomial kernel.
Take $q=32$, $R=655351$, and $r_0=(R-1)/2=327675$. Then
$\operatorname{popcount}(r_0)=16$, and for the kernel's exact half-binomial
ordinate $Y$ one has
\[
 v_2(Y)=15,\quad X=2^R\equiv2\pmod{31},\quad
 Y\equiv23\pmod{31},\quad H\equiv0\pmod{31}.
\]
Thus $q^3\mid Y$ but $\gcd(H,q-1)=31$. The modular residue is checked
by a finite Lucas-digit computation, cross-checked against direct
binomial sums on small inputs. This is a kernel example, not a zero for
the authentic fixed outer tuple. Its role is to show why the proof must
not assume universal coprimality.

The integer companion uses a related but simpler idea: normalize the
auxiliary equation to adjacent coefficients before merging its root.
The descent lemma restores integrality that a bare resultant would not
guarantee. Both improvements change the dependency order of the proof,
not merely the naming of an existing arithmetic register.

\section{Reproducibility, evidence, and remaining work}

\subsection{What is supplied and checked}
The archive contains the complete report source and PDF, the literal
positive and natural arithmetic graphs, the integer residual packets,
all required coefficient payloads, deterministic builders, independent
checkers, upstream proof notes, source pins, and attribution notices.

The ten-variable verification has distinct layers:
\begin{enumerate}
\item Reconstruct the entire coefficient tuple from the literal machine,
      verify its finite geometry and padding certificates, and compare every
      integer byte with the bound payloads.
\item Verify the four-root norm by exact polynomial arithmetic, independently
      of the circuit builder. No random evaluation stands in for that identity.
\item Check the graph's topology, liveness, full input set, exact schedule,
      operation count, degree bound, and nonzero leading certificate.
\item Compare independently implemented evaluations using the actual fixed
      coefficients modulo primes, as transcription diagnostics.
\item Exercise controls for radical cancellation, nonadjacent Pell descent,
      and the failure of a coprimality assumption. Mutating a critical norm
      coefficient must be rejected.
\end{enumerate}
The integer component separately checks monic restoration, centered-root
parity, two-root elimination, exact degrees, and its fixed coefficient
bindings. The root runner records each component's result.

\subsection{What finite verification cannot establish}
Finite tests do not prove the unbounded Pell or compiler theorems. Those
are mathematical arguments with named source dependencies. The report's
new arithmetic proof has received an independent adversarial review, but
it is not a kernel-checked Lean or Rocq theorem. The native compiler's
object-language ZFC internalization remains open in the reviewed snapshot.

The construction is explicit in the standard computational sense: a finite
graph and fully specified integer coefficients. Expanding every enormous
coefficient and every monomial into a single line would not simplify the
mathematics. The complete fixed integer bytes and rebuilding recipes are
included rather than replaced by an unspecified existence claim.

\subsection{Useful next steps}
The immediate verification target is the new residue-separation theorem,
with its sign recovery and canonical completeness, then the literal
four-root graph. The theorem is modular: it requires the same normalized
source interface already isolated in the repository.

A later size search can center the new carrier or reduce its majorant,
provided the rational positivity bootstrap survives. Another possible
direction is to adapt the same small-residue principle to different
proof-certificate encodings. Neither proposed extension is counted as a
completed equation in this report. The delivered results are the concrete
ten-positive-variable and integer companion constructions above.

\appendix
\section{Fixed coefficients and exact archive definitions}\label{app:coefficients}

For $P_{10}$, each listed gzip payload decompresses to a byte string
$b_0,\ldots,b_{m-1}$ and denotes exactly
$\sum_{j=0}^{m-1}b_j256^j$. There is no sign byte or truncation.
\begin{center}
\begin{tabular}{@{}lrrl@{}}
\toprule
Port & Bits & Raw bytes & Payload\\
\midrule
$\beta$ &590069&73759&\file{beta.bin.gz}\\
$K_0$ &592884&74111&\file{K.bin.gz}\\
$M_C$ &590069&73759&\file{MC.bin.gz}\\
$M_F$ &590070&73759&\file{MF_source.bin.gz}\\
\midrule
Total&&295388&\\
\bottomrule
\end{tabular}
\end{center}
The full hashes are in \file{ten_positive/coefficients/BINDING.json}.
The unchanged compact44 generator is retained with its attribution and
its literal transition data. \file{rebuild_coefficients.py} reconstructs
all four values in a temporary directory and compares the complete raw
integer bytes. The native radix is $\rho=2^{67}$ with cell length 8807;
therefore $B=\rho^{8807}=2^{590069}$. The compiler uses the literal
943-tile language, complement fields, duplicated clauses, and its
specified padding controller \cite{compact44}. All those finite constants
are fixed by the included recipe.

The integer companion's six coefficient ports are separately stored and
bound under \file{integer/coefficients/}. They are based on
$\beta=2^{1064921}-1$, with fixed cone scales $2^{1113}$ and $2^{666104}$.
The compact44 payloads must not be substituted into that construction.

\section{Files and commands}
From the archive root, run
\begin{quote}\ttfamily
python3 verify\_all.py
\end{quote}
The main ten-variable component uses only the Python standard library.
The integer symbolic builder has its pinned requirements in its own
\file{requirements.txt}; the archive's README explains installation when
needed. To rebuild the report, run
\begin{quote}\ttfamily
pdflatex -interaction=nonstopmode -halt-on-error report.tex\\
pdflatex -interaction=nonstopmode -halt-on-error report.tex
\end{quote}
No script needs a GitHub connection to define or evaluate the delivered
polynomials. All proof-source paths are pinned in source manifests.

\newpage
\begin{thebibliography}{99}
\raggedright
\bibitem{active} V. Reshetnikov, \emph{ProveIt: DiophantineConsistency active
research}, snapshot \pin, especially \file{README.md},
\file{SIZE_REDUCTION.md}, \file{STATUS.md}, and \file{OBJECT_ZFC.md}.
\url{https://github.com/VladimirReshetnikov/ProveIt/tree/c1a102936aca437e9999530a0111e6ddab56dc6d/Computability/DiophantineConsistency}.
\bibitem{compact44} ProveIt package 44, \emph{Complement-packed Wang943
coefficients}, \file{44-b4-sparse/native/PROOF.md}, coefficient recipe,
and bindings. Pinned sources and actual bytes retained in this archive.
\bibitem{carriers} ProveIt package 61, \emph{Congruence carriers},
\file{report/congruence_carriers.tex}, \file{positive/PROOF13.md}, and
\file{positive/PROOF14.md}. The normalized source interface, direct
strong-rank lemma, signed auxiliary-index lemma, and canonical witness
construction are used explicitly.
\bibitem{cones} ProveIt package 69, \emph{Weighted cones and rational
Pell descent}, especially \file{positive/PROOF.md}.
\bibitem{roottransfer} ProveIt package 70, \emph{Residue orientation and
root transfer}, especially the eleven-positive-variable proof and signed
power-recovery theorem.
\bibitem{joint} \emph{ZFC Diophantine Joint Guards}, incoming archive
\file{ZFC_Diophantine_Joint_Guards_2026-10-24.zip}, at the reviewed
snapshot. Used components: \file{pell/}, \file{joint/}, and
\file{circuits/}. The archive's UTC date differs from this report's
Washington State date.
\bibitem{structural} \emph{Structural Reductions, Round 11}, incoming
archive \file{ConZFC_Structural_Reductions_Round11_2026-10-24.zip},
especially the transport pivot and calibrated cone proof.
\bibitem{radical} \emph{Radical guard descent}, incoming archive
\file{con-zfc-radical-guard-descent-2026-10-23.zip}; used for the current
comparison frontier.
\bibitem{budget} \emph{Budget and centering}, incoming archive
\file{con-zfc-budget-and-centering-2026-10-23.zip}; in particular its
four-radical separation discussion and carrier-compatibility boundary.
\bibitem{kernel} ProveIt native stream queue, \emph{A strong half-binomial
Pell kernel in 42 operations}, \file{pell_kernel_half_binomial42.md}.
The retained note supplies the exact half-binomial ordinate and valuation
formula used only in the coprimality control.
\bibitem{sun} Z.-W. Sun, \emph{Ten unknowns for Hilbert's tenth problem
over the integers}, arXiv:2609.01594v2 (2026).
\url{https://arxiv.org/abs/2609.01594}.
\end{thebibliography}
\end{document}
